\section{Contextualization}

Distributed computing emerged as a vital paradigm to address the fundamental limitations of centralized architecture, particularly when faced with the exponential growth of Big Data, complex simulations, and modern artificial intelligence. For decades, vertical scaling (scale-up) sufficed as hardware manufacturers consistently increased single-core processor speeds in accordance with Moore's Law \parencite{hennessy2019architecture}. 

However, physical constraints such as quantum uncertainties and overheating and eventually created a severe infrastructure bottleneck, preventing Moore's Law's from staying true \parencite{waldrop2016moore}. By shifting the focus to horizontal scaling (scale-out), distributed computing connects an interconnected network of commodity computers, or nodes, to function as a unified system. This approach successfully mitigates the financial burden of high-end hardware, eliminates single points of failure through native fault tolerance, and solves the computational bottleneck by breaking down massive datasets into smaller, parallelizable tasks \parencite{sevilla2015framework}.

The historical trajectory of this field reflects a steady evolution from isolated mainframes to global, ubiquitous networks. The foundational concepts were established in the 1960s and 1970s with early packet-switching networks like ARPANET, which transitioned in the 1980s into localized clusters tied together by early message-passing interfaces. By the late 1990s, the paradigm expanded into Grid Computing, utilizing idle worldwide resources for grand-challenge scientific problems \parencite{foster2003grid}. 

The definitive turning point arrived in the 2000s when technology giants introduced frameworks like MapReduce, which later inspired open-source ecosystems such as Apache Hadoop and Spark. Coupled with the rise of modern cloud computing, distributed computing evolved from an exclusive academic tool into the essential backbone of contemporary digital infrastructure, enabling real-time, resilient, and scalable data analytics across the globe \parencite{dean2004mapreduce}.

However, the impact of this MapReduce-driven revolution expanded far beyond commercial data mining and web search algorithms. This fundamental capability to fragment massive datasets into parallel, independent tasks became the vital engine for accelerating high-fidelity complex simulations. Fields demanding intensive mathematical modeling such as particle physics, computational fluid dynamics and aerospace simulations historically collided with the stark limitations of sequential processing. By translating the divide-and-conquer concepts popularized by MapReduce into \gls{hpc} environments \parencite{ekanayake2010twister}, modern distributed systems can now simulate hyper-complex scenarios in record time \parencite{reed2015exascale}. What once required months of processing on a single centralized machine is now executed in a matter of hours or minutes across computing clusters, effectively transforming computing time from a scientific bottleneck into a catalyst for technological innovation.

\section{Motivation}

The amount of simulation data produced nowadays is also growing rapidly, especially in environments where multiple scenarios, parameter sweeps, and repeated executions are needed to support analytical decision making. In many scientific and industrial contexts, the main challenge is no longer only to execute simulations, but also to interpret the large volumes of output that they generate. Once the execution stage is complete, analysts still need to extract useful information from raw results, compare alternative scenarios, and transform those outputs into metrics that can support robust conclusions. This problem becomes particularly relevant when a simulation batch contains many runs and each run produces detailed records that must later be processed, filtered, and aggregated \parencite{li2018postprocessing,lagares2022aquila}.

This situation leads to two immediate concerns. The first is how to store large amounts of simulation output in a way that remains scalable and accessible. The second is how to process these outputs efficiently once they have been generated. Even when the analysis required for each individual simulation is conceptually simple, the total cost may become substantial when the number of runs is large. A single machine may become a bottleneck both in storage access and in serial processing time, especially if the analyst must first move the raw data to a local environment before any meaningful reduction occurs. For this reason, distributed storage and distributed computing appear not as optional refinements, but as natural responses to the scale of the problem \parencite{ghemawat2003gfs,dean2004mapreduce}.

Because of these challenges, modern computing has produced frameworks and tools that allow users to express large-scale distributed computations without manually orchestrating every communication and synchronization detail. Different systems emphasize different trade-offs, but most of them aim to provide the same core properties: scalability, fault tolerance, and a programming model that lets the user focus more on the analytical task than on the infrastructure itself. In the context of data-intensive analytics, the MapReduce model provided a foundational abstraction for batch-oriented distributed computation \parencite{dean2004mapreduce}, while later systems such as Spark expanded this model through richer direct acyclic graphs based execution strategies suited to iterative and structured workloads \parencite{zaharia2012rdd,armbrust2015sparksql}. In this sense, this work is not isolated from the broader history of distributed systems; rather, it is positioned within a long effort to make large-scale data processing both practical and reusable.

For simulation-based decision support, this discussion has a very concrete implication. If the outputs of a simulation batch can be partitioned, stored, and processed in parallel, then the time between executing the simulations and obtaining analytical insight can be reduced substantially. The analyst no longer needs to treat post-processing as a separate, mostly manual phase performed after data retrieval. Instead, the infrastructure itself can participate in the production of metrics and aggregate results. This is precisely the motivation of the present work: to understand how distributed-processing ideas can be applied not to the execution of the simulation itself, but to the \gls{aa} stage that transforms raw simulation outputs into decision-support products.

\section{Objective}

This work aims to design and implement a distributed computing architecture for aerospace simulation environments. This architecture will be called ''AeroSOuDA'': \textbf{Aero}space \textbf{S}imulation \textbf{Ou}tput \textbf{D}istributed \textbf{A}rchitecture. The goal is to integrate the final solution to \gls{asa}, the project that motivated this work in the first place. In order achieve this goal, some rules must be established:

\begin{itemize}
    \item The architecture must support a workflow in which the analyst can define or select via interface a metric function and an aggregate function before simulation batch execution;
    \item The architecture must support any kind of metric and aggregate function;
    \item The architecture must have a \gls{dfs} to store simulation results;
    \item The primary language used will be Python, considering the robust ecosystem it brings, allowing us to analyse frameworks such as PySpark;
    \item The architecture must be tested in a local computer cluster beforehand;
    \item Considering a simulation batch, simulation results may an come at any order;
    \item Each metric processing should run preferably in the same node its corresponding simulation was processed in order to reduce data transit via network;
    \item The aggregate result shall assume at least 2 output types: numbers or graphic plots.
\end{itemize}

 These elements together define the intended contribution of the work: not a new simulation engine, but a new way to transform simulation outputs into useful analytical results more efficiently. By contrast, any improvement toward the simulation processing workflow is outside of the scope. The work does not aim to redesign the simulation engine, change the runtime logic by which individual simulations are executed, or optimize the internal physics or tactical behavior of the simulator. 
